\section{Height data and the central Chow projector}\label{sec:setup}

\subsection{Height data and integral images}\label{sec:height-data}

Fix an algebraic closure \(\Fbar\) of \(F\) and an embedding into
\(\C\). Homological triviality is defined by geometric rational
Betti cohomology. All Chow groups below have rational coefficients
unless otherwise indicated. We regard \(V^r(X_E)\) as a subspace
of \(V^r(X_{\Fbar})\): restriction is injective because a rational
equivalence descends to a finite extension and trace composed with
restriction multiplies by its degree.

\begin{definition}\label[definition]{def:heights}
A curve-normalized height system consists of pairings
\[
 h_X:V^r(X_E)\times V^{n+1-r}(X_E)\longrightarrow\R
\]
over every finite extension \(E/F\), satisfying:
\begin{enumerate}[label=(H\arabic*)]
\item bilinearity over \(\Q\) and symmetry under exchange of the
complementary degrees;
\item invariance under extension in the absolute normalization,
and covariance under isomorphisms, including Galois conjugation;
\item correspondence adjunction
\[
 h_X(\Gamma_*z,w)=h_Y(z,{}^t\Gamma_*w);
\]
\item on a pointed curve \(C\), the form \(-h_C\) on
\(\Pic^0(C)_\Q\) is the positive Néron--Tate form for the canonical
principal polarization, in the corresponding Poincaré normalization.
\end{enumerate}
\end{definition}

These hypotheses will be understood in every statement involving
\(h\). A real positive form means positive definiteness after
tensoring the rational group over a fixed number field with \(\R\).
For codimension \(r\) on \(X\), the symbol \(\AJ_\Q\)
denotes the complex Abel--Jacobi map with target
\(J^r(X_\C)(\C)\otimes_\Z\Q\).

Let \(\Lambda_X^r(E)\) be the image in \(V^r(X_{\Fbar})\) of the
integral-coefficient cycles over \(E\) with zero geometric rational
cycle class. Thus torsion and rational equivalence are removed before
counting. For a fixed rational projector \(P\), we use its projected
integral image \(P_*\Lambda_X^r(E)\); for an integral projector this
equals \(\Lambda_X^r(E)\cap P_*V^r(X_E)\).

\begin{definition}\label[definition]{def:northcott}
For compatible integral images \(\Lambda(E)\) in one geometric
rational group and a quadratic form \(q\), bounded-degree Northcott
means that, for every \(B\in\R\) and integer \(D\ge1\),
\[
 \{\xi\in\bigcup_{[E:F]\le D}\Lambda(E):q(\xi)\le B\}
\]
is finite. Fixed-field Northcott means the same assertion with
\(E\) fixed. Bounded denominators mean the union of
\(d^{-1}\Lambda(E)\) for \(1\le d\le d_0\), with \(d_0\) fixed.
\end{definition}

Northcott for integral images implies the bounded-denominator version
by scaling. A fixed initial extension \(F_0/F\) is harmless for
bounded-degree assertions: after replacing \(E\) by \(EF_0\),
\([EF_0:F_0]\le[E:F]\), the integral image remains integral, and
the absolute height is unchanged.

\subsection{Curve-invisible cycles}\label{sec:invisible-cycles}

\begin{lemma}\label[lemma]{lem:invisible}
Let \(C\) be a smooth projective curve and
\(\Gamma\in\CH^r(C\times X)_\Q\), with \(\dim X=n\). If
\(\Gamma_*H^1(C)=0\) in \(H^{2r-1}(X)\), then
\[
 {}^t\Gamma_*w=0\quad\text{in }\Pic^0(C)_\Q
 \qquad(w\in V^{n+1-r}(X)).
\]
Hence \(\Gamma_*\Pic^0(C)_\Q\) is in the left height radical.
\end{lemma}

\begin{proof}
By Poincaré duality, the transpose map from
\(H^{2n+1-2r}(X)\) to \(H^1(C)\) is zero. Functoriality of the
Abel--Jacobi map under Chow correspondences therefore makes the
Abel--Jacobi class of \({}^t\Gamma_*w\) zero. On a curve, the
rationalized Abel--Jacobi map is the identification with
\(\Pic^0(C)_\Q\). Vanishing after complex base change descends for
divisor classes with rational coefficients. The final assertion is
adjunction. The correspondence functoriality used here is also
proved in \cite[Lemma~4.11]{LagaShnidman}.
\end{proof}

Algebraically trivial cycles are generated by images of degree-zero
divisors on smooth projective curves. For geometric algebraic
triviality over \(E\), the finitely many data can be defined over a
finite extension and then traced to \(E\). This convention is used
throughout.

\subsection{The maximal curve-visible projector}
\label{sec:central-projector}

Let \(k=2c-1\), \(Y=\prod_{i=1}^kC_i\), and
\(\Pi=\bigotimes_i\pi_{1,i}\). Set
\[
 U_E=\Pi_*V^c(Y_E),\qquad
 A_E=U_E\cap\CH^c(Y_E)_{\alg,\Q},\qquad G_E=U_E/A_E.
\]
For the product polarization \(\ell_Y=\sum a_i[o_i]\), the actual
Chow identity \(\ell_Y\Pi=0\) follows from
\([o_i]\pi_{1,i}=0\). Thus the geometric cohomology of \(\Pi\)
is middle and primitive.

Let \(W\subset\Pi H^k(Y,\Q)\) be the span of all images
\[
 H^1(C)\longrightarrow\Pi H^k(Y,\Q)
\]
of codimension-\(c\) curve correspondences over finite extensions
of \(F\). Choose finitely many correspondences spanning \(W\),
and one field \(F_0\) defining them and their source points.
Write their matrix as \(\Gamma=(\Gamma_j)_j\), with
\(\Gamma_j=\Pi\Gamma_j\pi_{1,j}\), and put
\(N=\bigoplus_jh^1(C_j)\), \(J=\prod_j\Pic^0(C_j)\).

\begin{proposition}\label[proposition]{prop:central-projector}
There are actual Chow projectors \(e,R\) over \(F_0\), with
\[
 e={}^te,\quad e^2=e,\quad e\Pi=\Pi e=e,\quad R=\Pi-e,
\]
such that the following hold for every finite \(E/F_0\).
\begin{enumerate}[label=(\roman*)]
\item The cohomological action of \(e\) is the orthogonal
projection onto \(W\). There is a fixed abelian variety \(B/F_0\)
with \(eU_E\simeq B(E)_\Q\), through fixed Chow correspondences.
The form \(q_e=(-1)^ch_Y\) on this summand is a positive
Néron--Tate form.
\item With \(K_E=RA_E\),
\[
 K_E=\ker(\AJ_\Q|_{A_E})
    =A_E\cap\rad(h_Y|_{U_E})
    =\rad(h_Y|_{A_E}).
\]
\item There is an orthogonal decomposition
\(U_E=eU_E\oplus RU_E\). The residual form on
\(RU_E/RA_E\simeq G_E\) is independent of the choice of maximal
curve-visible projector. Write it as \(h_G\) and set
\(q_G=(-1)^ch_G\).
\end{enumerate}
\end{proposition}

\begin{proof}
Set \(a={}^t\Gamma\Gamma\). Since \(2c-k=1\), its entries are
divisor correspondences, so \(a\in\End(N)=\End(J)_\Q\).
This realization is faithful, and transposition corresponds to
Rosati adjunction for the product principal polarization.
Geometric Hodge--Riemann on the primitive subspace \(W\) implies
that \((-1)^{c-1}a\) is positive semidefinite and
\(\ker a=\ker\Gamma\) on \(H^1\). In particular \(a\) is
semisimple. Choose a polynomial \(f\in\Q[t]\) equal to
\(t^{-1}\) on the nonzero spectrum and zero at zero, and set
\[
 a^+=f(a),\qquad p=a^+a,\qquad
 e=\Gamma a^+{}^t\Gamma.
\]
Faithfulness on \(N\) makes \(a^+aa^+=a^+\), \(p^2=p={}^tp\)
actual Chow identities. These identities prove the assertions about
\(e\) and its cohomology. Let \(B=pJ\), up to isogeny.
The maps \(\Gamma\) and \(a^+{}^t\Gamma\) identify
\(B(E)_\Q\) with \(eU_E\). Adjunction gives
\[
 h_Y(\Gamma x,\Gamma y)=h_{\mathrm{curves}}(x,ay).
\]
The curve normalization and the positivity of
\((-1)^{c-1}a\) on \(pJ\) give (i).

For every later curve correspondence \(\Delta\), the geometric
\(H^1\) action of \(R\Delta\) is zero by maximality of \(W\).
Curve generation and \Cref{lem:invisible} give
\(RA_E\subseteq\rad(h_Y|_{U_E})\). The Abel--Jacobi map is zero
on \(RA_E\) by the same geometric action. Its restriction to
\(eU_E\) is injective after rationalization: the corresponding
map from \(B\) to the intermediate Jacobian has injective rational
\(H_1\), hence finite kernel. Since \(eU_E\subseteq A_E\) and
its height is nondegenerate, all equalities in (ii) follow.

The self-transpose projectors give the orthogonal decomposition.
The radical assertion just proved makes the residual form well
defined on \(G_E\). Two orthogonalized lifts of the same Griffiths
class differ by an element of \(A_E\) orthogonal to \(A_E\).
By (ii), their difference is in the full radical. Thus the residual
form is independent of the lift and of \(e\).
\end{proof}

The cohomological subspace \(W\) and its orthogonal complement are
defined independently of the chosen spanning correspondences. Betti
and étale comparison transports them to each \(\Q_{\ell_0}\)
realization. The algebraic construction of \(e\) in the proof uses
geometric Hodge--Riemann and the faithful realization of \(h^1\)
endomorphisms; it is available independently of the height data.

\begin{lemma}\label[lemma]{lem:visible-integral}
The projected integral image in \(eU_E\) has bounded-degree
Northcott for \(q_e\), allowing any fixed denominator bound.
\end{lemma}

\begin{proof}
Put \(D=a^+{}^t\Gamma\). These fixed curve correspondences satisfy
\(De=D\). One integer \(N\) clears their denominators, so an
integral cycle \(z\) gives a genuine point \(NDz\in J(E)\).
The same point detects \(ez\), even though \(ez\) may have
rational coefficients.

On the fixed factor \(pJ(\Fbar)_\Q\), the ambient Néron--Tate
height of \(J\) is bounded above by a fixed positive multiple of
\(q_e\) under the inverse \(D\). This is comparison of two
polarizations on \(B\); passing to rational points removes
torsion without changing either height. Thus a bound for
\(q_e(ez)\) bounds the height of the actual point \(NDz\).
Classical bounded-degree Northcott on \(J\), followed by the
injectivity of \(D\) on \(eU_E\), proves the assertion. Fixed
extra denominators are absorbed into \(N\). See
\cite[Sections~9--10]{ConradFeng} for the height and Mordell--Weil
facts used here.
\end{proof}
